\documentclass[conference]{IEEEtran}
\IEEEoverridecommandlockouts
\usepackage{cite}
\usepackage{amsmath,amssymb,amsfonts}
\usepackage{algorithmic}
\usepackage{graphicx}
\usepackage{textcomp}
\usepackage{xcolor}
\usepackage{booktabs}
\usepackage{multirow}
\usepackage{array}
\usepackage{url}
\usepackage{hyperref}
\usepackage{balance}

\hypersetup{
  colorlinks=true,
  linkcolor=black,
  citecolor=black,
  urlcolor=black
}

\graphicspath{{paper_assets/figures/}}

\def\BibTeX{{\rm B\kern-.05em{\sc i\kern-.025em b}\kern-.08em
    T\kern-.1667em\lower.7ex\hbox{E}\kern-.125emX}}

\begin{document}

\title{Beyond Controlled-Domain Accuracy: A Leakage-Aware Crop Disease
Classification Pipeline with Calibration, Selective Prediction,
and Multilingual Decision Support}

\author{
\IEEEauthorblockN{Teja Sai Maddirala}
\IEEEauthorblockA{\textit{School of Computer Science and Engineering} \\
\textit{VIT-AP University}\\
Amaravati, India \\
teja.23bce7573@vitapstudent.ac.in}
\and
\IEEEauthorblockN{Prajwal Aaryan Immadi}
\IEEEauthorblockA{\textit{School of Computer Science and Engineering} \\
\textit{VIT-AP University}\\
Amaravati, India \\
prajwal.23bce8631@vitapstudent.ac.in}
\linebreakand
\IEEEauthorblockN{Sai Nihar Kunala}
\IEEEauthorblockA{\textit{School of Computer Science and Engineering} \\
\textit{VIT-AP University}\\
Amaravati, India \\
nihar.23bce7213@vitapstudent.ac.in}
%
\and
%
\IEEEauthorblockN{Tanishq S}
\IEEEauthorblockA{\textit{School of Computer Science and Engineering} \\
\textit{VIT-AP University}\\
Amaravati, India \\
tanishq.23bce20211@vitapstudent.ac.in}
\linebreakand
\IEEEauthorblockN{Koduru Hajarathaiah\textsuperscript{*}}
\IEEEauthorblockA{\textit{School of Computer Science and Engineering} \\
\textit{VIT-AP University}\\
Amaravati, India \\
hajarathaiah.k@vitap.ac.in}
}
\maketitle

\begin{abstract}
Crop disease classification systems built on controlled laboratory datasets
frequently exhibit near-perfect internal accuracy yet falter when exposed to
field-acquired imagery. This paper presents a reproducible, leakage-aware
pipeline for 19-class crop disease recognition derived from a curated
PlantVillage subset spanning five crops: bell pepper, corn, potato, soybean,
and tomato. The pipeline combines rigorous duplicate auditing and locked
three-way data partitioning with convolutional and transformer classifiers,
post-hoc temperature-scaling calibration, confidence-based selective prediction,
a lesion-grounded heuristic severity proxy, zero-shot external evaluation on
PlantDoc, and a retrieval-augmented generation (RAG) multilingual advisory layer.
EfficientNet-B0 trained on 18,113 images achieves 99.53\% top-1 accuracy
and 0.9929 macro-F1 on the locked 5,322-image internal test set. Temperature
scaling reduces the expected calibration error (ECE) from 0.00321 to 0.00228
without altering classification decisions. An EfficientNetV2-S proposed model
and a ViT-B/16 comparator reach comparable internal accuracy but offer no
material advantage over the baseline. Zero-shot evaluation on a 144-image,
15-class semantically mapped subset of PlantDoc reveals sharp domain shift:
top-1 accuracy drops to 19.44\% and macro-F1 to 0.1428, while top-5 accuracy
remains 67.36\%, indicating that the correct class is retrievable within the top
predictions but is not reliably top-ranked under field-like conditions. These
findings demonstrate that high controlled-domain accuracy is insufficient
evidence for field deployment readiness. The RAG advisory layer, validated
across all 19 labels in English, Hindi, and Telugu, provides safety-bounded
educational guidance conditional on classifier confidence, explicitly prohibiting
pesticide-prescription outputs and mandating expert confirmation below
high-confidence thresholds.
\end{abstract}

\begin{IEEEkeywords}
crop disease detection, EfficientNet, ViT, temperature scaling, calibration,
selective prediction, domain shift, PlantDoc, retrieval-augmented generation,
multilingual advisory, PlantVillage
\end{IEEEkeywords}

% =====================================================================
\section{Introduction}
% =====================================================================

Crop diseases are responsible for substantial yield losses worldwide, with
estimates from the Food and Agriculture Organization of the United Nations (FAO)
suggesting that plant pathogens and pests account for up to 40\% of annual
global crop production losses~\cite{savary2019global,fao2021state}.
For smallholder farmers with limited access to extension services, early and
accurate disease identification is a critical determinant of treatment timing
and economic outcome~\cite{pound2017uncertainty,bhatt2022farmer}. Automated
computer-vision systems trained on large annotated datasets offer the potential
to democratise field-level disease diagnosis~\cite{mohanty2016deep,ferentinos2018deep},
particularly when deployed on mobile platforms accessible in low-connectivity
regions~\cite{ramcharan2017deep,atila2021plant}.

The PlantVillage dataset~\cite{hughes2015open,sladojevic2016deep}, comprising
over 54,000 images collected under controlled, single-leaf conditions against
uniform backgrounds, has become the dominant benchmark for crop disease
classification. Models trained on PlantVillage routinely achieve accuracies
exceeding 95\% and often surpass 99\%~\cite{mohanty2016deep,too2019comparative,kamilaris2018deep}.
However, a growing body of evidence cautions against interpreting such figures
as deployment readiness. Controlled-domain datasets exhibit a fundamental
mismatch with the complex backgrounds, variable lighting, overlapping leaves,
and partial occlusions encountered in real agricultural
fields~\cite{ferentinos2018deep,thenmozhi2019crop,chen2020using}.
This domain shift problem has been demonstrated empirically through the
PlantDoc dataset~\cite{singh2020plantdoc}, which provides field-acquired images
and enables direct measurement of the gap between controlled and field-like
performance.

Beyond raw accuracy, well-calibrated confidence estimates are essential for safe
model deployment. Overconfident predictions---wherein the model assigns high
softmax probabilities to incorrect classes---are particularly hazardous in
high-stakes agricultural decisions~\cite{guo2017calibration,minderer2021revisiting}.
Temperature scaling~\cite{guo2017calibration} is a widely used post-hoc
calibration method that rescales logits by a learned scalar, improving
confidence alignment without altering predictions. Selective prediction
frameworks~\cite{geifman2017selective,el2010foundations} complement calibration
by allowing the model to abstain on low-confidence inputs, providing a mechanism
for reducing error rates at the cost of partial coverage.

The explainability of neural-network predictions has also attracted considerable
attention in the agricultural domain~\cite{selvaraju2017gradcam,chattopadhay2018gradcam++,
lundberg2017shap}. Lesion-localisation approaches can highlight regions of
classifier attention and, with appropriate caveats, provide heuristic visual
severity estimates~\cite{saleem2019plant, albattah2022novel}.

Finally, even when a reliable visual classifier is available, farmers require
actionable guidance in their native language~\cite{sujatha2021performance}.
Retrieval-augmented generation (RAG) systems~\cite{lewis2020retrieval,guu2020retrieval}
can serve as a structured advisory interface that grounds natural-language
responses in vetted agronomic knowledge bases while enforcing safety constraints
against prescriptive pesticide recommendations.

\subsection*{Research Objectives}

This study addresses the following objectives:
\begin{enumerate}
    \item Construct a reproducible, leakage-free 19-class crop disease classification
          pipeline on PlantVillage-derived data using strict partition isolation.
    \item Evaluate and compare EfficientNet-B0, EfficientNetV2-S, and ViT-B/16
          classifiers on controlled-domain internal metrics, calibration quality,
          and selective prediction behaviour.
    \item Quantify zero-shot domain shift by evaluating the best-performing model
          on the external PlantDoc test set under a transparent mapping protocol.
    \item Introduce a lesion-grounded heuristic severity proxy with explicitly
          stated limitations as a qualitative grounding tool.
    \item Design and validate a safety-bounded multilingual RAG advisory layer
          that conditions guidance on classifier confidence bands.
\end{enumerate}

\subsection*{Contributions}

The principal contributions are:
\begin{itemize}
    \item A fully reproducible, near-duplicate-audited 19-class data split
          protocol with documented seed, class mapping, and partition counts.
    \item Rigorous multi-architecture benchmarking on the same locked partitions,
          showing that a compact EfficientNet-B0 matches heavier models at
          substantially lower latency and superior calibration.
    \item The first transparent PlantDoc zero-shot external evaluation of this
          pipeline, quantifying the controlled-to-field performance gap.
    \item A qualitative lesion-grounded severity proxy with documented protocol
          limitations, explicitly precluding agronomic or clinical severity claims.
    \item A prototype multilingual RAG advisory layer validated across all 19
          labels in English, Hindi, and Telugu, with a tiered confidence-safety policy.
\end{itemize}

% =====================================================================
\section{Related Work}
% =====================================================================

\subsection{Plant Disease Datasets and Benchmarking}

Hughes and Salathé~\cite{hughes2015open} introduced PlantVillage as the first
large-scale open-access leaf-disease image repository. Subsequent work has
critiqued the controlled capture conditions as unrepresentative of field
settings~\cite{mohanty2016deep,sladojevic2016deep,barbedo2018factors,barbedo2019plant}.
Singh et al.~\cite{singh2020plantdoc} addressed this gap with PlantDoc, a
2,600-image multi-crop dataset collected under natural field conditions.
PlantDoc exposes the performance drop experienced by models trained exclusively
on PlantVillage-style imagery. Additional datasets targeting specific crops,
including cassava~\cite{mwebaze2019icassava,ramcharan2017deep}, rice~\cite{prajapati2017analysis},
wheat~\cite{zhang2020crop}, and banana~\cite{magnier2021detection} have been
developed, but PlantVillage remains the standard multi-crop controlled benchmark.
Dataset governance issues such as duplicate images across splits have been
identified as a source of inflated accuracy estimates, motivating careful
de-duplication~\cite{cheng2021fabricated,northcutt2021pervasive,barz2020we}.

\subsection{CNN-Based Plant Disease Classification}

Mohanty et al.~\cite{mohanty2016deep} pioneered the application of deep
convolutional neural networks (CNNs) to PlantVillage, achieving 99.35\% accuracy
under colour-augmented training. Ferentinos~\cite{ferentinos2018deep} reported
similarly high accuracy using AlexNet, VGG, and GoogLeNet variants. The effect
of transfer learning from ImageNet-pretrained weights has been consistently
observed to accelerate convergence and improve final accuracy on small-to-medium
agricultural datasets~\cite{too2019comparative,kamilaris2018deep,tm2018detection}.
Densely connected architectures such as DenseNet~\cite{huang2017densely} and
residual networks~\cite{he2016deep} have been widely adopted in disease-focused
applications~\cite{agarwal2020neural,Chen2020identify}. Data augmentation
strategies including random cropping, horizontal flipping, brightness jitter, and
colour-space transformations are standard practice to improve
generalisation~\cite{shorten2019survey,cubuk2020randaugment}.

\subsection{Efficient Neural Architectures}

Tan and Le introduced EfficientNet~\cite{tan2019efficientnet}, which applies
compound scaling to simultaneously balance network width, depth, and input
resolution. EfficientNets achieve state-of-the-art accuracy with substantially
fewer parameters than preceding architectures. The successor
EfficientNetV2~\cite{tan2021efficientnetv2} further improved training speed
through progressive learning and Fused-MBConv blocks. Applications in
agricultural disease recognition have demonstrated that EfficientNet variants
offer a favourable accuracy-efficiency tradeoff~\cite{ahmed2022automated,
howlader2022effects,mahum2023novel}. MobileNet~\cite{howard2017mobilenets}
and its successors~\cite{sandler2018mobilenetv2,howard2019searching} provide
even lighter alternatives for on-device inference at the cost of accuracy.

\subsection{Vision Transformers in Agricultural Applications}

Dosovitskiy et al.~\cite{dosovitskiy2021image} introduced the Vision Transformer
(ViT) that processes image patches as token sequences with self-attention. Data
Efficient Image Transformers (DeiT)~\cite{touvron2021training} addressed ViT's
large data requirement through distillation. ViT and Swin
Transformer~\cite{liu2021swin} have been applied to plant disease recognition,
sometimes matching or exceeding CNN performance on benchmark
datasets~\cite{chen2022vit,alruwaili2023efficient,yu2023leaf}. However,
transformer architectures typically require more computation per image, limiting
their practical advantage in resource-constrained deployments. Hybrid CNN-ViT
models~\cite{wu2021cvt,chen2021crossvit} aim to combine local feature extraction
with global context modelling.

\subsection{Probability Calibration and Selective Prediction}

Calibration refers to the agreement between a model's predicted confidence and
its empirical accuracy~\cite{guo2017calibration}. Temperature scaling is a
one-parameter post-hoc calibration method that divides logits by a scalar $T$
optimised to minimise negative log-likelihood on a held-out calibration
set~\cite{guo2017calibration,platt1999probabilistic}. Because top-1 predictions
are invariant to monotone confidence transformations, temperature scaling changes
ECE~\cite{naeini2015obtaining} and Brier score~\cite{brier1950verification} but
not accuracy or macro-F1. Selective classification~\cite{geifman2017selective,
el2010foundations,franc2023optimal} introduces an abstention option: the model
declines to classify inputs below a chosen confidence threshold, trading coverage
for improved accuracy on the retained subset. Evaluating the risk--coverage
tradeoff~\cite{geifman2019selectivenet} provides a more complete picture of
deployment safety than a single accuracy figure. Overconfidence under distribution
shift has been highlighted as a critical failure mode for safety-critical
applications~\cite{ovadia2019can,lakshminarayanan2017simple}.

\subsection{Domain Adaptation and External Validation}

The domain gap between controlled-laboratory and field-acquired data is a
well-documented challenge in agricultural computer vision~\cite{barbedo2018factors,
sun2016deep}. Approaches for bridging this gap include domain adaptation~\cite{ganin2015unsupervised,long2015learning},
self-supervised pre-training on unlabelled field images~\cite{chen2020simple,he2020momentum},
and test-time augmentation~\cite{krizhevsky2012imagenet}. However, direct
zero-shot evaluation---applying a source-domain-trained model to target-domain
data without any adaptation---remains the most conservative and honest estimate
of model robustness under distribution shift~\cite{torralba2011unbiased}.
Several studies have demonstrated large performance drops when PlantVillage
models are evaluated on field data~\cite{singh2020plantdoc,chen2020using,
thenmozhi2019crop,fuentes2017robust}, underscoring the importance of external
benchmarking as part of model evaluation.

\subsection{Explainability, Lesion Analysis, and Severity Estimation}

Class Activation Mapping (CAM)~\cite{zhou2016learning} and Grad-CAM~\cite{selvaraju2017gradcam}
are widely used to visualise the regions of an image that most influence
classification decisions. LIME~\cite{ribeiro2016should} and SHAP~\cite{lundberg2017shap}
provide model-agnostic attribution. Lesion segmentation methods range from
threshold-based colour separation~\cite{phadikar2013rice} to deep semantic
segmentation~\cite{long2015fully,ronneberger2015unet}. Severity estimation---
the quantification of disease extent within a leaf---typically requires expert-annotated
pixel masks for validation~\cite{bock2010plant,francl1997quantifying}. In the
absence of validated annotations, any automatic severity estimate must be
characterised as a heuristic
proxy~\cite{mahlein2012development,barbedo2014review}.

\subsection{Retrieval-Augmented Generation for Agricultural Advisory}

RAG systems~\cite{lewis2020retrieval} combine a retrieval module over a curated
knowledge base with a language model to generate grounded, source-supported
responses. Applications in the agricultural domain~\cite{gupta2023indiaai,
meena2023plantgpt,xu2023agri} demonstrate the utility of RAG for providing
context-sensitive pest management guidance. Responsible deployment of AI advisory
systems in agriculture requires explicit safety constraints, uncertainty
communication, and avoidance of unverified pesticide or chemical
recommendations~\cite{rolnick2022tackling,jobin2019global,
floridi2019establishing}. Multilingual capability is particularly important
for serving smallholder farmers in regions such as South Asia~\cite{sujatha2021performance,
dwivedi2022multilingual}.

% =====================================================================
\section{Materials and Methods}
% =====================================================================

\subsection{Datasets, Class Definition, and Data Governance}

The primary dataset is a curated subset of PlantVillage~\cite{hughes2015open},
selecting five economically important crop species: bell pepper
(\textit{Capsicum annuum}), corn (\textit{Zea mays}), potato (\textit{Solanum
tuberosum}), soybean (\textit{Glycine max}), and tomato (\textit{Solanum
lycopersicum}). The resulting pool contained 31,729 images across 20
class--crop combinations. Five images identified as exact validation-set
duplicates were removed prior to any split, leaving 31,724 clean images.

The primary classification task, designated \textbf{Main19}, uses 19 classes,
excluding \texttt{soybean\_\_healthy}, which was reserved for a secondary
auxiliary experiment (Aux20). The 19 Main19 classes span eight disease conditions
and five healthy states across the four represented crops (corn, potato, bell
pepper, tomato), reflecting the 13 diseased and six healthy categories available
after applying the exclusion. The largest class (tomato yellow leaf curl virus)
contains 5,357 images; the smallest (potato healthy) contains 152 images,
yielding a class imbalance ratio of 35.2. Fig.~\ref{fig:split_distribution}
shows the image counts per class across all three partitions.

The external evaluation dataset is PlantDoc~\cite{singh2020plantdoc}, a
field-acquired, multi-crop leaf-disease dataset. Only the PlantDoc test split
was used; training images from PlantDoc were not accessed or used by any model
in this study for training, model selection, temperature fitting, or threshold
calibration. This ensures that the PlantDoc evaluation is strictly zero-shot.

\begin{figure}[t]
\centering
\includegraphics[width=\columnwidth]{fig_01_split_distribution.pdf}
\caption{Image counts per Main19 class across the three locked partitions
(training: 18,113; calibration: 3,197; internal test: 5,322). Classes are
ordered by total size. All partitions were stratified by class and their
membership frozen before any model training.}
\label{fig:split_distribution}
\end{figure}

\subsection{Duplicate Auditing and Reproducible Data Splits}

Near-duplicate detection was conducted using perceptual hashing (pHash) with a
Hamming-distance threshold of $\leq$4 bits, yielding 136 near-duplicate
candidate pairs~\cite{zauner2010implementation}. Each candidate pair was
subjected to manual visual review; two image pairs were confirmed as true
cross-split near-duplicates and removed, preventing potential
leakage~\cite{barz2020we,northcutt2021pervasive}. All downstream experiments
were conducted on the resulting de-duplicated image pool.

Reproducible three-way stratified splits were constructed at random seed 7.
The original PlantVillage training pool served as the development pool from
which 15\% were stratified-sampled to form the calibration split. The original
PlantVillage validation set was retained unchanged as the locked internal test
set. Partition membership was frozen and saved as signed manifests before any
model training commenced, ensuring no post-hoc modification could occur.

The final locked Main19 partition counts are:
\textbf{train}: 18,113 images across 19 classes;
\textbf{calibration}: 3,197 images;
\textbf{internal test}: 5,322 images.

\subsection{Image Classification Models and Training Protocol}

Three architectures were evaluated on the Main19 task. All models were
initialised with ImageNet-pretrained weights (IMAGENET1K\_V1) and fine-tuned
on the training split. Checkpoint selection was governed by calibration-set
macro-F1. The internal test set was accessed only once, for final one-shot
evaluation after checkpoint selection.

\textbf{EfficientNet-B0 (Baseline).}
The baseline classifier uses EfficientNet-B0~\cite{tan2019efficientnet}
with a 19-class linear head. Images were resized to $224\times224$ pixels.
Training used class-weighted cross-entropy loss with AdamW
optimiser~\cite{loshchilov2017decoupled} at learning rate $3\times10^{-4}$,
weight decay $10^{-4}$, batch size 32, and a maximum of 15 epochs with early
stopping after 4 non-improving calibration epochs. Step-decay halved the
learning rate at epoch 6 and again at epoch 13. The final checkpoint was
selected at epoch 15 (calibration macro-F1: 0.9940).

\textbf{Proposed EfficientNetV2-S.}
The proposed model replaces the backbone with EfficientNetV2-S~\cite{tan2021efficientnetv2},
a larger and more recent architecture optimised for training efficiency. Input
resolution was increased to $384\times384$ pixels. Training used class-weighted
focal cross-entropy loss~\cite{lin2017focal} ($\gamma=1.5$) with AdamW at
learning rate $2\times10^{-4}$, weight decay $10^{-4}$, batch size 12, and a
maximum of 20 epochs (early stopping patience: 5). Checkpoint selection again
used calibration macro-F1. The best checkpoint was selected at epoch 13
(calibration macro-F1: 0.9979).

\textbf{ViT-B/16 (Architecture Comparator).}
Vision Transformer ViT-B/16~\cite{dosovitskiy2021image} was trained as an
architecture comparator under the same protocol constraints. Input size was
$224\times224$. Training used class-weighted cross-entropy with label smoothing
($\varepsilon=0.05$)~\cite{szegedy2016rethinking} and AdamW at learning rate
$10^{-4}$, weight decay $10^{-4}$, batch size 16, and a maximum of 15 epochs
(early stopping patience: 4). The best checkpoint was selected at epoch 11
(calibration macro-F1: 0.9947).

All experiments were conducted in PyTorch 2.10.0+cu128, Python 3.12.13,
NumPy 2.0.2, and pandas 2.3.3 at seed 7. Training was performed on
NVIDIA GPU hardware via the Kaggle platform.

Table~\ref{tab:architecture} summarises the model configurations.

\begin{table}[t]
\caption{Model Architectures, Training Configuration, and Efficiency}
\label{tab:architecture}
\centering
\setlength{\tabcolsep}{4pt}
\begin{tabular}{lccc}
\toprule
\textbf{Property} & \textbf{EfficientNet-B0} & \textbf{EfficientNetV2-S} & \textbf{ViT-B/16} \\
\midrule
Input size & $224^2$ & $384^2$ & $224^2$ \\
Parameters & $\approx$5.3 M & $\approx$20.2 M & 85.8 M \\
Loss & Wtd.\ CE & Wtd.\ Focal CE & Wtd.\ CE + LS \\
Focal $\gamma$ & --- & 1.5 & --- \\
Label smooth. & No & No & 0.05 \\
Optimiser & AdamW & AdamW & AdamW \\
LR & $3\!\times\!10^{-4}$ & $2\!\times\!10^{-4}$ & $1\!\times\!10^{-4}$ \\
Weight decay & $10^{-4}$ & $10^{-4}$ & $10^{-4}$ \\
Batch size & 32 & 12 & 16 \\
Max epochs & 15 & 20 & 15 \\
ES patience & 4 & 5 & 4 \\
Sel.\ epoch & 15 & 13 & 11 \\
Latency (ms/img) & $\dagger$ & $\dagger$ & 17.4 \\
\bottomrule
\multicolumn{4}{l}{\small $\dagger$ Not measured; \small Wtd.\ CE = class-weighted cross-entropy;} \\
\multicolumn{4}{l}{\small LS = label smoothing; ES = early stopping.} \\
\end{tabular}
\end{table}

\subsection{Post-hoc Calibration and Selective Prediction}

\textbf{Temperature Scaling.}
Following the calibration framework of Guo et al.~\cite{guo2017calibration},
temperature scaling was applied to the EfficientNet-B0 baseline by fitting a
scalar temperature parameter $T$ that minimises negative log-likelihood on the
calibration split. The optimal temperature for the baseline was $T=1.5695$.
Because temperature scaling applies a monotone transformation to predicted
probabilities, it does not alter the top-1 class assignment, preserving
accuracy and macro-F1 identically. Its effect is measured by ECE, NLL, and
Brier score.

Expected calibration error (ECE)~\cite{naeini2015obtaining} was computed with
15 equal-width bins. Brier score~\cite{brier1950verification} measures the
mean squared error between predicted probability vectors and one-hot targets.

\textbf{Selective Prediction.}
Selective prediction was implemented for all three architectures following the
risk-coverage framework of Geifman and El-Yaniv~\cite{geifman2017selective}.
A confidence threshold $\tau$ was selected on the calibration split under the
constraint that the calibration-set error rate (risk) was $\leq 1\%$ while
maximising coverage. Thresholds were only ever selected on calibration data;
the internal test set was not used for threshold tuning. For all three models,
the selected threshold was $\tau = 0.50$, which also maximised coverage.
The selected threshold was then applied unchanged to the internal test set as
a purely prospective evaluation.

\subsection{Lesion Grounding and Severity Proxy}

A heuristic lesion-grounded severity proxy was computed for all images in the
Main19 dataset. The proxy estimates the fraction of a leaf's pixel area that
exhibits colour or texture patterns associated with visible lesions, applying
background separation followed by HSV/Lab abnormal-colour and low-green
vegetation heuristics. The method requires no separate lesion segmentation model
and no expert annotation.

The proxy assigns each image to one of three ordinal bands:
\textit{mild\_proxy} ($<5\%$ visible lesion area),
\textit{moderate\_proxy} ($5$--$20\%$), and
\textit{severe\_proxy} ($>20\%$). These labels are visual-heuristic designations
only. The protocol metadata explicitly states that the proxy: (i) is not
expert-annotated lesion segmentation; (ii) is not a clinical disease-stage label;
(iii) has not been validated against agronomic treatment thresholds; and
(iv) does not account for hidden lesions, leaf age, whole-plant burden, or
temporal progression. All severity claims in this paper are accordingly qualified
as proxy estimates.

The proxy was applied to the full Main19 pool (20,562 diseased-class images after
excluding healthy classes from the lesion computation). Results are reported by
class and by split.

\subsection{Zero-Shot External-Domain Evaluation}

To measure the practical impact of source-to-target domain shift, the
EfficientNet-B0 baseline---the best-performing model under calibration-set
selection---was applied zero-shot to the PlantDoc test split. ``Zero-shot'' in
this context means no PlantDoc image was used for training, model selection,
calibration, or threshold tuning.

PlantDoc organises images into 27 folder categories. Of these, 15 could be
unambiguously mapped to Main19 classes. The remaining 12 folders cover crops
absent from Main19 (apple, blueberry, cherry, peach, raspberry, soybean, squash,
strawberry, grape) and were excluded. Mapped classes and image counts are
detailed in Table~\ref{tab:plantdoc_mapping}. The mapping was constructed from
folder names and verified against disease taxonomy; mapping ambiguity is
explicitly discussed as a limitation.

Evaluation was conducted on the model's unscaled softmax outputs (temperature
$T=1.0$); temperature scaling was not applied because the internal calibration
temperature was fitted on PlantVillage calibration data, not PlantDoc data,
making its application to PlantDoc an unjustified extrapolation.
Evaluation metrics include top-1 accuracy, top-5 accuracy, balanced accuracy,
macro-F1, and weighted F1. Confidence distributions for correct and incorrect
predictions were examined separately.

\begin{table*}[t]
\caption{PlantDoc-to-Main19 Class Mapping and Zero-Shot Evaluation Inventory.
Twelve PlantDoc folders covering crops absent from Main19 were excluded.}
\label{tab:plantdoc_mapping}
\centering
\begin{tabular}{llrr}
\toprule
\textbf{PlantDoc Folder} & \textbf{Mapped Main19 Label} & \textbf{Images} & \textbf{Top-1 Acc.} \\
\midrule
Bell\_pepper\_leaf & bell\_pepper\_\_healthy & 8 & 0.750 \\
Bell\_pepper\_leaf\_spot & bell\_pepper\_\_bacterial\_spot & 9 & 0.000 \\
Corn\_Gray\_leaf\_spot & corn\_\_cercospora\_leaf\_spot\_gray\_leaf\_spot & 4 & 0.500 \\
Corn\_rust\_leaf & corn\_\_common\_rust & 10 & 0.200 \\
Corn\_leaf\_blight & corn\_\_northern\_leaf\_blight & 12 & 0.250 \\
Potato\_leaf\_early\_blight & potato\_\_early\_blight & 14 & 0.286 \\
Potato\_leaf\_late\_blight & potato\_\_late\_blight & 8 & 0.000 \\
Tomato\_leaf\_bacterial\_spot & tomato\_\_bacterial\_spot & 9 & 0.000 \\
Tomato\_Early\_blight\_leaf & tomato\_\_early\_blight & 9 & 0.111 \\
Tomato\_leaf & tomato\_\_healthy & 8 & 0.000 \\
Tomato\_leaf\_late\_blight & tomato\_\_late\_blight & 10 & 0.600 \\
Tomato\_mold\_leaf & tomato\_\_leaf\_mold & 6 & 0.000 \\
Tomato\_Septoria\_leaf\_spot & tomato\_\_septoria\_leaf\_spot & 12 & 0.167 \\
Tomato\_leaf\_mosaic\_virus & tomato\_\_tomato\_mosaic\_virus & 10 & 0.000 \\
Tomato\_leaf\_yellow\_virus & tomato\_\_tomato\_yellow\_leaf\_curl\_virus & 15 & 0.133 \\
\midrule
\multicolumn{2}{l}{\textbf{Total compatible}} & \textbf{144} & \textbf{0.194} \\
\bottomrule
\end{tabular}
\end{table*}

\subsection{Retrieval-Grounded Multilingual Advisory Layer}

The RAG advisory layer provides safety-bounded educational guidance following
a predicted diagnosis. It is not an image classifier and does not alter
classification outputs. The knowledge base contains 19 entries, one per Main19
label, with structured agronomic content including symptom descriptions, causal
agents, management principles, and prevention guidelines. Content was authored
and curated specifically for this study; no external unverified content was
incorporated.

Retrieval uses TF-IDF cosine similarity with classifier-label
prioritisation~\cite{robertson2009bm25}: when a classifier label is provided,
that entry is retrieved directly and supplemented by secondary similarity search.
Responses are generated in three languages: English, Hindi, and Telugu~\cite{bird2009nlp}.

The advisory layer enforces a tiered confidence-safety policy:
\textit{low confidence} ($<0.70$): requests clearer images and expert
confirmation; \textit{medium confidence} ($0.70$--$0.89$): provides
screening-level guidance and recommends inspecting multiple plants;
\textit{high confidence} ($\geq 0.90$): provides substantive agronomic guidance
with mandatory field-verification warnings. Across all tiers, the system
explicitly prohibits pesticide brand, dosage, or location-specific chemical
recommendations, and it states explicitly that its output is educational and
does not constitute a confirmed agronomic or clinical diagnosis.

The RAG layer was validated by running all 19 Main19 labels through the retrieval
pipeline and verifying that: (i) the primary retrieved label matched the input
label; and (ii) an advisory of appropriate length and content was generated.
All 19 retrievals were successful (100\% primary label match), as shown in
Table~\ref{tab:rag_validation}.

% =====================================================================
\section{Results}
% =====================================================================

\subsection{Dataset Characteristics and Split Integrity}

The Main19 split yielded 18,113 training, 3,197 calibration, and 5,322 internal
test images. Stratification preserved the class proportions across all three
partitions, as confirmed by Fig.~\ref{fig:split_distribution}. The class
imbalance ratio (35.2) is primarily driven by the large tomato yellow leaf curl
virus class (5,357 total) and the small potato healthy class (152 total). Class
weights were applied during training to compensate for imbalance.

Near-duplicate auditing confirmed two cross-split near-duplicate pairs from the
initial 136 pHash candidates. Both images were removed before split creation.
Five exact duplicates had been removed even earlier during initial EDA. No
further cross-partition leakage was detected. The manifests were locked and no
partition membership was changed after training commenced.

\subsection{Controlled-Domain Classification Performance}

\textbf{EfficientNet-B0 Baseline.}
The baseline model achieves 99.53\% top-1 accuracy, 99.32\% balanced accuracy,
and 0.9929 macro-F1 on the locked internal test set. Fig.~\ref{fig:learning_curves}
shows the training and calibration loss and macro-F1 across 15 epochs. Training
converged steadily; calibration macro-F1 reached 0.9940 at the final epoch,
which was selected as the best checkpoint.

\begin{figure}[t]
\centering
\includegraphics[width=\columnwidth]{fig_02_baseline_learning_curves.pdf}
\caption{EfficientNet-B0 training and calibration macro-F1 and loss across 15
epochs. Checkpoint selection was made at epoch 15, corresponding to the highest
calibration macro-F1 (0.9940). The internal test set was not accessed during
training.}
\label{fig:learning_curves}
\end{figure}

\textbf{EfficientNetV2-S Proposed Model.}
EfficientNetV2-S achieves 99.81\% accuracy and 0.9978 macro-F1 on the
calibration split (best checkpoint at epoch 13). On the internal test set,
as evaluated in the multi-model comparison (Notebook 07), EfficientNetV2-S
reaches 99.55\% accuracy and 0.9926 macro-F1, matching but not exceeding the
baseline when evaluated on the same held-out partition. Temperature-scaled
internal test metrics for EfficientNetV2-S were not independently extracted
and are reported as NOT VERIFIED for the test partition.

\textbf{ViT-B/16.}
ViT-B/16 achieves 99.36\% accuracy and 0.9926 macro-F1 on the internal test set
(best checkpoint at epoch 11). Its calibration macro-F1 was 0.9947.

Table~\ref{tab:internal_comparison} presents the consolidated multi-architecture
comparison. EfficientNet-B0 and EfficientNetV2-S reach nearly identical
accuracy (99.55\% versus 99.55\% in the shared comparison table), while
ViT-B/16 is marginally lower (99.36\%). NLL is substantially higher for
ViT-B/16 (0.1481 vs.\ 0.0159 for B0), indicating poorer confidence calibration
before scaling. Brier score confirms this ordering:
B0 records 0.000337, EfficientNetV2-S 0.00773, and ViT-B/16 0.0311.

\begin{table}[t]
\caption{Controlled-Domain Internal Test Performance Comparison (Main19,
5,322 images). B0 = EfficientNet-B0; V2-S = EfficientNetV2-S. Higher is better
for Acc., Bal.\ Acc., and Macro-F1; lower is better for NLL and Brier.}
\label{tab:internal_comparison}
\centering
\setlength{\tabcolsep}{3pt}
\begin{tabular}{lrrrrr}
\toprule
\textbf{Model} & \textbf{Acc.} & \textbf{Bal.\ Acc.} & \textbf{Macro-F1} & \textbf{NLL} & \textbf{Brier} \\
\midrule
B0 & 0.9955 & 0.9935 & 0.9931 & 0.0159 & 0.000337 \\
V2-S & 0.9955 & 0.9935 & 0.9926 & 0.0176 & 0.00773 \\
ViT-B/16 & 0.9936 & 0.9938 & 0.9926 & 0.1481 & 0.0311 \\
\bottomrule
\multicolumn{6}{l}{\small Source: \texttt{07.../cnn\_vs\_vit\_internal\_test\_comparison.csv}.} \\
\end{tabular}
\end{table}

\textbf{Model Selection.}
EfficientNet-B0 is identified as the primary model for calibration, external
evaluation, and RAG integration based on: (i) lowest Brier score and NLL on the
internal test set, indicating better calibration even without explicit scaling;
(ii) matching or exceeding V2-S and ViT accuracy at substantially lower
parameter count ($\approx$5.3 M vs.\ $\approx$20.2 M and 85.8 M respectively);
and (iii) absence of confirmed temperature-scaled test metrics for V2-S, making
full comparison infeasible. ViT-B/16 is retained as an architecture comparator
only; its substantially higher latency (17.4 ms/img) and poorer calibration
make it a less suitable primary deployment choice.

\subsection{Calibration and Confidence-Aware Prediction}

\textbf{Temperature Scaling (Baseline).}
The optimal temperature for EfficientNet-B0 was $T=1.5695$, fitted on the
calibration split. As expected, temperature scaling does not alter the top-1
classification decision: accuracy (0.9953) and macro-F1 (0.9929) on the internal
test set are identical before and after scaling. However, calibration metrics
improve substantially. Table~\ref{tab:calibration_summary} presents the full
calibration profile.

\begin{table}[t]
\caption{EfficientNet-B0 Calibration Metrics Before and After Temperature
Scaling ($T=1.5695$). Accuracy and Macro-F1 are invariant to scaling.
ECE computed with 15 equal-width bins. Lower is better for NLL, Brier, ECE.}
\label{tab:calibration_summary}
\centering
\setlength{\tabcolsep}{4pt}
\begin{tabular}{llrrr}
\toprule
\textbf{Split} & \textbf{Scaling} & \textbf{NLL} & \textbf{Brier} & \textbf{ECE} \\
\midrule
\multirow{2}{*}{Calibration} & Uncalibrated & 0.01850 & 0.00699 & 0.00392 \\
 & Temp.\ scaled & 0.01830 & 0.00678 & 0.00130 \\
\midrule
\multirow{2}{*}{Internal test} & Uncalibrated & 0.01550 & 0.00704 & 0.00321 \\
 & Temp.\ scaled & 0.01593 & 0.00642 & \textbf{0.00228} \\
\bottomrule
\multicolumn{5}{l}{\small Source: \texttt{04.../main19\_temperature\_scaling\_metrics.csv},} \\
\multicolumn{5}{l}{\small \texttt{04.../test\_evaluation\_run\_metadata.json}.} \\
\end{tabular}
\end{table}

ECE decreases from 0.00321 to 0.00228 on the internal test set (a 29\% relative
reduction). Brier score improves from 0.00704 to 0.00642. The calibration
reliability diagrams confirm this improvement: the temperature-scaled model
exhibits substantially closer alignment between mean confidence bins and empirical
accuracy, as illustrated in Fig.~\ref{fig:reliability}.

\begin{figure}[t]
\centering
\includegraphics[width=\columnwidth]{fig_03_test_reliability_calibrated.pdf}
\caption{EfficientNet-B0 calibration reliability diagram on the internal test
set after temperature scaling ($T=1.5695$). Each bin shows the mean predicted
confidence (horizontal) versus empirical accuracy (diagonal parity). The
temperature-scaled model closely tracks the identity diagonal, indicating
well-calibrated confidence estimates. ECE: 0.00228 (15 bins).}
\label{fig:reliability}
\end{figure}

\textbf{Selective Prediction.}
For all three models, the confidence threshold selected on the calibration split
to achieve $\leq1\%$ calibration risk was $\tau=0.50$. At this threshold,
EfficientNet-B0 accepted all 5,322 internal test images (coverage: 100\%),
achieving a test selective accuracy of 99.55\% and risk of 0.0045.
EfficientNetV2-S rejected 2 images (coverage: 99.96\%), ViT-B/16 rejected 3
(coverage: 99.94\%). The calibration-derived threshold transferred reliably to
the internal test set across all three architectures, confirming the validity of
the calibration-only threshold selection protocol.

Fig.~\ref{fig:risk_coverage} shows the risk--coverage curves on the internal
test set. At the selected threshold, all three models operate near full coverage
with risk below 0.5\%. EfficientNet-B0 and ViT-B/16 achieve full-test accuracy
of 0.9953--0.9955 under the threshold; all models show monotone risk reduction
as coverage decreases with higher thresholds, as expected. Table~\ref{tab:selective}
summarises the selective prediction results.

\begin{figure}[t]
\centering
\includegraphics[width=\columnwidth]{fig_04_risk_coverage_curves.pdf}
\caption{Risk--coverage curves on the internal test set for all three classifiers.
Risk (top-1 error rate on accepted images) decreases monotonically as the
confidence threshold rises and coverage decreases. The vertical dashed line
marks the selected threshold $\tau=0.50$, chosen on the calibration set only.}
\label{fig:risk_coverage}
\end{figure}

\begin{table}[t]
\caption{Selective Prediction Results on the Internal Test Set.
Threshold $\tau=0.50$ was selected on the calibration set. Acc.\ (full) is
accuracy on all test images regardless of threshold.}
\label{tab:selective}
\centering
\setlength{\tabcolsep}{4pt}
\begin{tabular}{lrrrrr}
\toprule
\textbf{Model} & $\bm{\tau}$ & \textbf{Cov.} & \textbf{Reject.} & \textbf{Sel.\ Acc.} & \textbf{Risk} \\
\midrule
EfficientNet-B0 & 0.50 & 1.000 & 0 & 0.9955 & 0.0045 \\
EfficientNetV2-S & 0.50 & 0.9996 & 2 & 0.9957 & 0.0043 \\
ViT-B/16 & 0.50 & 0.9994 & 3 & 0.9955 & 0.0045 \\
\bottomrule
\multicolumn{6}{l}{\small Source: \texttt{08.../selective\_prediction\_test\_results.csv}.} \\
\end{tabular}
\end{table}

\subsection{Lesion-Grounded Severity-Proxy Findings}

The heuristic lesion proxy was computed for all diseased-class Main19 images
across training, calibration, and internal test splits. Fig.~\ref{fig:lesion}
presents a qualitative panel of sample images across severity bands.

\begin{figure}[t]
\centering
\includegraphics[width=\columnwidth]{fig_05_lesion_severity_proxy.pdf}
\caption{Qualitative panel of lesion-grounded severity proxy estimates for
selected disease classes. Colour-highlighted regions indicate pixels identified
as lesion-like by the heuristic masking procedure. Proxy bands are designated
\textit{mild\_proxy} ($<$5\% lesion area), \textit{moderate\_proxy} (5--20\%),
and \textit{severe\_proxy} ($>$20\%). These are visual-heuristic designations;
no expert annotation or agronomic validation was used.}
\label{fig:lesion}
\end{figure}

Across all splits, the mean visible lesion area is approximately 6.2\%
(calibration: 6.19\%; internal test: 6.22\%; training: 6.26\%), with a
median of approximately 4.4\% and a standard deviation of approximately 6.5\%.
The close agreement across splits indicates that partition stratification
preserved lesion-area distributions effectively.

Notable class-level variation is observed. Classes with large visible lesions
include potato early blight (mean 12.5\%), corn gray leaf spot (13.0\%),
and potato late blight (11.0\%). Classes with smaller lesion areas include
tomato mosaic virus (1.8\%), corn healthy (0.9\%), and tomato leaf mold
(3.4\%). These differences are consistent with the known visual appearance of
each disease: early blight and late blight produce large necrotic lesions,
whereas mosaic virus primarily causes leaf curling and colour mottling without
prominent discrete lesion areas.

The proxy values are not validated against expert-annotated lesion ground truth.
They provide exploratory characterisation and a qualitative visual grounding of
classifier attention regions. They must not be interpreted as clinically or
agronomically validated severity estimates.

\subsection{External Domain Shift on PlantDoc}

The EfficientNet-B0 baseline was applied zero-shot to the 144-image,
15-class PlantDoc external test subset. The results reveal severe domain shift.

\textbf{Overall metrics.}
Top-1 accuracy: 0.1944 (28 of 144 correct). Balanced accuracy: 0.1998.
Macro-F1: 0.1428. Weighted F1: 0.1698. Top-5 accuracy: 0.6736.

The contrast with internal performance (99.53\% top-1 accuracy) is stark and
expected: the model was trained on controlled, single-leaf images against
uniform backgrounds, while PlantDoc contains field-acquired images with complex
natural backgrounds, variable lighting, and co-occurring foliage.

\textbf{Top-5 analysis.}
The high top-5 accuracy (67.4\%) indicates that the correct class is retrievable
within the model's top-5 predictions for the majority of images, even when it
is not the top-ranked class. This suggests that the model retains some ability
to distinguish relevant disease features under domain shift, but that ranking
confusion---driven by visually similar diseases under different imaging
conditions---prevents reliable top-1 identification.

\textbf{Confidence analysis.}
Despite poor accuracy, the model remains overconfident: the mean confidence on
all 144 PlantDoc images is 0.845, with a median of 0.942. Correct predictions
have a mean confidence of 0.883; incorrect predictions have a mean confidence
of 0.836. The small gap between correct and incorrect confidence distributions
(0.048) indicates that confidence alone is insufficient to identify domain-shifted
misclassifications, unlike the internal test setting where high-confidence errors
are rare. Importantly, even at the internal-test-derived threshold of $\tau=0.50$,
only 11 of 144 images are rejected (see Table~\ref{tab:plantdoc_selective}),
and accepted image accuracy improves only marginally to 19.5\%, confirming that
the calibration-derived threshold provides minimal protection against domain-shift
errors.

\begin{table}[t]
\caption{PlantDoc Zero-Shot External Evaluation.
Top-5 accuracy indicates the correct label appears within the top-5 predictions.
Temperature $T=1.0$ (unscaled). Confidence from calibration-derived
$\tau=0.50$ selective threshold provides only marginal improvement.}
\label{tab:plantdoc_selective}
\centering
\setlength{\tabcolsep}{4pt}
\begin{tabular}{lrrrrr}
\toprule
\textbf{Model} & \textbf{Imgs} & \textbf{Top-1} & \textbf{Top-5} & \textbf{Macro-F1} & \textbf{Mean Conf.} \\
\midrule
EfficientNet-B0 & 144 & 0.1944 & 0.6736 & 0.1428 & 0.845 \\
\bottomrule
\multicolumn{6}{l}{\small Source: \texttt{10.../table\_05\_external\_domain\_overall\_metrics.csv},} \\
\multicolumn{6}{l}{\small \texttt{10.../run\_metadata.json}. Classes: 15 compatible Main19 classes.} \\
\end{tabular}
\end{table}

\textbf{Per-class and per-crop analysis.}
Class-level results (Table~\ref{tab:plantdoc_mapping}) show that six of 15
classes achieve zero correct top-1 predictions: bell pepper bacterial spot,
potato late blight, tomato bacterial spot, tomato healthy, tomato leaf mold,
and tomato mosaic virus. The highest-performing classes are bell pepper healthy
(F1: 0.24) and tomato late blight (F1: 0.28). At the crop level, bell pepper
achieves the highest top-1 accuracy (0.35), followed by corn (0.27), potato
(0.18), and tomato (0.14).

Fig.~\ref{fig:plantdoc_confusion} shows the external-domain confusion matrix,
revealing systematic misclassification patterns. Fig.~\ref{fig:plantdoc_crop}
shows top-1 accuracy by crop.

\begin{figure}[t]
\centering
\includegraphics[width=\columnwidth]{fig_06_plantdoc_confusion_matrix.pdf}
\caption{Confusion matrix for EfficientNet-B0 on the PlantDoc external test set
(144 images, 15 compatible classes). Rows: true labels; columns: predicted
labels. The low diagonal density confirms severe top-1 accuracy degradation
under domain shift. Many images are misclassified into high-prevalence training
classes.}
\label{fig:plantdoc_confusion}
\end{figure}

\begin{figure}[t]
\centering
\includegraphics[width=\columnwidth]{fig_07_plantdoc_accuracy_by_crop.pdf}
\caption{PlantDoc zero-shot top-1 accuracy broken down by crop species. Bell
pepper achieves the highest accuracy (0.35); tomato the lowest (0.14). Image
counts per crop are annotated. Differences reflect both the number of compatible
classes and the severity of domain shift per crop.}
\label{fig:plantdoc_crop}
\end{figure}

\subsection{RAG Advisory Evaluation}

Table~\ref{tab:rag_validation} presents the validation summary for the RAG
advisory layer. All 19 Main19 labels retrieved the correct primary advisory
entry (100\% label-match recall). Advisory responses ranged from 1,568 to
1,914 characters, covering symptom descriptions, causal agents, and management
principles at the medium-confidence tier. The system correctly enforced the
confidence-safety policy: medium-confidence inputs generated screening-level
advisory text; high-confidence inputs generated substantive guidance with
mandatory field-verification warnings; low-confidence inputs requested clearer
imagery and expert confirmation.

\begin{table}[t]
\caption{RAG Advisory Validation Results.
Retrieval Method: TF-IDF cosine similarity. Languages: English, Hindi,
Telugu. All 19 Main19 labels tested at medium confidence (0.70--0.89 band).
Primary label match: 19/19 (100\%).}
\label{tab:rag_validation}
\centering
\begin{tabular}{lcc}
\toprule
\textbf{Metric} & \textbf{Value} & \textbf{Source} \\
\midrule
Knowledge base entries & 19 & Run metadata \\
Primary label match & 19/19 (100\%) & Table 01 \\
Languages supported & 3 (EN, HI, TE) & Run metadata \\
Advisory length range (chars) & 1,568--1,914 & Table 01 \\
Pesticide prescription prohibited & Yes & Safety policy \\
Expert confirmation required ($<0.70$) & Yes & Safety policy \\
Field-verification warning ($\geq 0.90$) & Yes & Safety policy \\
\bottomrule
\multicolumn{3}{l}{\small Source: \texttt{11.../table\_01\_rag\_validation.csv},} \\
\multicolumn{3}{l}{\small \texttt{11.../table\_02\_confidence\_safety\_policy.csv}.} \\
\end{tabular}
\end{table}

The RAG layer is a prototype educational advisory component. Its evaluation
is limited to structural validation: correct retrieval and policy enforcement.
It does not constitute an end-to-end user study, clinical validation, or
randomised controlled evaluation of agronomic guidance quality. These are
explicitly noted as directions for future work.

% =====================================================================
\section{Discussion}
% =====================================================================

\subsection{The Controlled-Domain Accuracy Ceiling and Its Limits}

The EfficientNet-B0 baseline achieves 99.53\% top-1 accuracy on the locked
internal test set---a result consistent with the broader literature on PlantVillage
models~\cite{mohanty2016deep,too2019comparative,ferentinos2018deep}. However,
this study deliberately places this result in context: it is a ceiling observed
under matched imaging conditions, not a guarantee of field-level performance.
The deliberate inclusion of calibration analysis, selective prediction, and
external evaluation provides a multi-dimensional view that raw accuracy alone
cannot offer.

The near-identical performance of EfficientNetV2-S and ViT-B/16 on the internal
test set, despite their substantially larger parameter counts and training
budgets, indicates that the controlled PlantVillage task is saturated---the
discriminative features required for near-perfect classification are accessible
to all three architectures. In this regime, the differentiating factors become
calibration quality, inference efficiency, and generalisation under shift.
EfficientNet-B0 wins on all three: it exhibits the lowest Brier score
(0.000337 versus 0.00773 and 0.0311), matches accuracy, and is estimated to
be substantially faster than ViT-B/16 (17.4 ms/img), making it the most
practically appropriate choice for the studied conditions.

\subsection{Calibration as a Prerequisite for Safe Deployment}

Temperature scaling with $T=1.5695$ reduces ECE on the internal test set from
0.00321 to 0.00228, a 29\% reduction with no cost to discrimination. This
improvement matters most in the selective prediction context: confidence
thresholds used for abstention decisions must be meaningful, i.e., higher
confidence must correspond to higher empirical accuracy. The reliability curves
confirm this property after scaling.

However, the PlantDoc confidence analysis reveals a critical limitation of
calibration under distribution shift: despite being calibrated on PlantVillage
data, the model assigns high confidence (median 0.942) to most PlantDoc images,
including a large majority of misclassified ones. The mean confidence gap between
correct and incorrect external predictions (0.048) is negligible compared to
the internal test regime where errors are rare and high-confidence errors are
identifiable. This finding confirms the observation in the literature that
post-hoc calibration on source-domain data does not generalise to
target-domain confidence estimates~\cite{ovadia2019can,guo2017calibration}.
A separate PlantDoc calibration set---which was not available and not used---
would be required to achieve meaningful external calibration.

\subsection{The External Domain Shift Gap}

The 80.1 percentage-point drop in top-1 accuracy between the internal test
(99.53\%) and the PlantDoc external test (19.44\%) is the most important finding
of this study. It should be interpreted as evidence of substantial source-to-target
domain shift rather than as a direct estimate of field-deployment performance
across all crops, regions, cultivars, and imaging conditions. The PlantDoc subset
used here comprises only 144 images from 15 classes, and several classes contain
as few as four images (corn gray leaf spot). Conclusions about specific class
difficulty must therefore be made cautiously.

Nevertheless, the result is consistent with published external evaluations of
PlantVillage-trained models on PlantDoc~\cite{singh2020plantdoc,chen2020using}
and reinforces that the controlled-domain accuracy advantage does not transfer
reliably. The top-5 accuracy of 67.4\% suggests that the model has not lost all
discriminative knowledge: the correct answer is within the model's top five
candidates for most images. This provides a basis for designing human-in-the-loop
systems that present multiple candidate diagnoses rather than forcing a single
output.

\subsection{Lesion Proxy and Explainability Limitations}

The heuristic lesion proxy provides a lightweight visual grounding mechanism
that does not require a separate segmentation model. The class-level patterns
are broadly consistent with known visual pathology---early blight and late blight
produce larger, more discrete lesions than mosaic virus or leaf curl. However,
the proxy has not been validated against expert-annotated lesion masks, and its
heuristic colour/texture rules may confound image artefacts, natural leaf
variation, or reflection with pathological tissue. The proxy should be treated
as an exploratory tool for qualitative analysis and future annotation
prioritisation, not as a disease severity estimator.

\subsection{RAG Advisory and Responsible Deployment}

The RAG advisory layer addresses the gap between classification output and
actionable field guidance. By conditioning advisory content on classifier
confidence and explicitly enforcing a no-pesticide-prescription policy, the
system avoids the most dangerous overclaims possible from an automated diagnosis
system. The three-language support (English, Hindi, Telugu) begins to address
the multilingual needs of the primary intended user base~\cite{sujatha2021performance}.

However, the advisory layer has significant limitations. The knowledge base
was curated for this study and has not been reviewed or endorsed by professional
plant pathologists. TF-IDF retrieval is sensitive to vocabulary and lacks
semantic understanding of ambiguous or partial queries. The safety policy
operates at fixed confidence thresholds that may not generalise to novel
imaging conditions. A rigorous evaluation would require expert review of advisory
quality, field user studies, and adversarial testing with out-of-domain inputs.

\subsection{Limitations}

The following limitations must be acknowledged:

\textit{Controlled source data.} PlantVillage images were collected under
laboratory conditions. The demonstrated external performance (19.44\% top-1 on
PlantDoc) confirms that this choice fundamentally limits generalisation.

\textit{PlantDoc mapping uncertainty.} Mapping PlantDoc folder names to Main19
classes required semantic interpretation; the \texttt{Tomato\_leaf} folder was
assumed to represent healthy tomato leaves. Mapping errors would inflate or
deflate individual class metrics.

\textit{PlantDoc subset size.} With only 144 images across 15 classes (mean
9.6 per class), class-level PlantDoc metrics have high variance and should be
interpreted with caution.

\textit{Missing EfficientNetV2-S calibration metrics.} Temperature-scaled
internal-test metrics for EfficientNetV2-S were not independently extracted;
only calibration-set metrics are confirmed. The calibration comparison is
therefore based only on the baseline model.

\textit{Proxy severity validation.} The lesion-area proxy has not been validated
against expert-annotated ground truth. All severity band assignments are
heuristic only.

\textit{RAG evaluation depth.} The RAG advisory layer was validated only for
structural correctness (label retrieval). Agronomic quality, user comprehension,
and safety in adversarial scenarios have not been evaluated.

\textit{No statistical significance tests.} No confidence intervals or
significance tests were estimated for reported metrics. Differences between
closely performing models (B0 vs.\ V2-S) should not be interpreted as
statistically meaningful without such analysis.

% =====================================================================
\section{Conclusion}
% =====================================================================

This paper presents a reproducible, leakage-aware crop disease classification
pipeline evaluated across controlled and external domains. The principal finding
is that high controlled-domain accuracy---here 99.53\% for EfficientNet-B0 on
Main19---does not transfer to field-like imagery: zero-shot evaluation on PlantDoc
yields 19.44\% top-1 accuracy despite a top-5 accuracy of 67.4\%. This gap must
be confronted honestly rather than deferred to future work, and it is a core
contribution of this study rather than merely a limitation.

Post-hoc temperature scaling ($T=1.5695$) improves ECE from 0.00321 to 0.00228
on the internal test set without altering predictions. Selective prediction with
a calibration-derived threshold of 0.50 achieves near-full coverage at sub-0.5\%
risk internally, but provides only marginal protection against domain-shift
errors externally. The lesion-grounded severity proxy reveals interpretable
class-level variation in visible lesion area consistent with known pathological
features, while its limitations as an unvalidated heuristic are explicitly
documented.

The multilingual RAG advisory layer bridges the gap between classification output
and usable field guidance, enforcing a tiered confidence-safety policy that
prohibits pesticide prescriptions and mandates expert confirmation below
high-confidence thresholds. It is validated as a structural prototype; clinical
and agronomic quality evaluation remains future work.

Future directions include: domain adaptation pre-training on unlabelled field
imagery; collection of expert-validated lesion annotations to upgrade the
severity proxy; expansion of the PlantDoc evaluation to more crops and regions;
integration with mobile deployment platforms; and rigorous field user studies of
the RAG advisory system with professional agronomists.

% =====================================================================
\section*{Declarations}
% =====================================================================

\textit{Data availability.} The PlantVillage dataset is publicly available
at \url{https://plantvillage.psu.edu/}. PlantDoc is available at
\url{https://github.com/pratikkayal/PlantDoc-Dataset}. The derived split
manifests, class maps, and evaluation metadata used in this study are
available in the project repository.

\textit{Code availability.} All analysis notebooks and reproducibility
artefacts are available in the project repository at random seed 7,
with Python 3.12.13 / PyTorch 2.10.0 / torchvision 0.25.0.

\textit{Conflicts of interest.} The authors declare no conflicts of interest.

\textit{Responsible-use statement.} The advisory layer described in this paper
is designed for educational and screening purposes only. It does not constitute
a confirmed agronomic or clinical diagnosis and must not be used to prescribe
pesticides, chemicals, or treatments without professional agronomic review.

\textit{Acknowledgements.} The authors acknowledge the use of Kaggle GPU
resources for all model training experiments.

\balance

\bibliographystyle{IEEEtran}
\bibliography{references}

\end{document}
